Predictive maintenance is a natural test bed for adaptive internal computation because many failure decisions depend on gradually accumulating evidence rather than a single snapshot. This paper presents OuroMaintain, a looped latent reasoning model for equipment health assessment and maintenance triage. The system combines a GRU window encoder, a shared recurrent latent block, and a confidence-based early-exit policy that allocates additional internal steps to ambiguous telemetry windows. The evaluation is structured as a benchmark matrix spanning C-MAPSS engine telemetry, IMS run-to-failure bearings, a reduced Paderborn machinery benchmark, and a reduced LBNL HVAC benchmark. Across every matched loop-versus-LLM comparison in that matrix, the adaptive looped model outperforms both generic pretrained and task-specific DistilBERT baselines trained on serialized telemetry summaries. On the most detailed comparison case, C-MAPSS FD001, the adaptive loop achieves the best test macro F1 score of 0.918 while using an average depth of only 1.21 steps and 0.17 ms/sample latency, compared with 0.889 macro F1 for the fixed-depth loop, 0.861 for the direct baseline, 0.318 for the generic pretrained LLM baseline, and 0.345 for the task-specific LLM at 32.11 ms/sample. On the reduced Paderborn split, the adaptive model improves validation macro F1 from 0.184 for both LLM baselines to 0.281 while preserving sub-millisecond sample latency, and on the reduced HVAC split it reaches 0.600 macro F1 as a transfer result. These results suggest that adaptive looping can improve maintenance classification quality across multiple domains while remaining far more compute efficient than both generic pretrained and task-specific text-serialized LLM alternatives.
